\subsection{ISOMORPHISMS AND TRANSPORT OF STRUCTURES}

Let \(\Sigma\) be a species of structures in a theory \(\mathfrak{C}\) , on \(n\) principal base sets \(x_{1},\ldots ,x_{n}\) , with \(m\) auxiliary base sets \(\mathbf{A}_1,\dots ,\mathbf{A}_m\) . Let \(\mathbf{S}\) be the echelon construction scheme on \(n + m\) letters which features in the typical characterization of \(\Sigma\) , and let \(\mathbf{R}\) be the axiom of \(\Sigma\) . In a theory \(\mathfrak{C}'\) which is stronger than \(\mathfrak{C}\) , let \(\mathbf{U}\) be a structure of species \(\Sigma\) on sets \(\mathbf{E}_1,\ldots ,\mathbf{E}_n\) (as principal base sets) and let \(\mathbf{U}'\) be a structure of the same species on sets \(\mathbf{E}_1',\ldots ,\mathbf{E}_n'\) . Finally, let \(f_{i}\) (in \(\mathfrak{C}'\) ) be a bijection of \(\mathbf{E}_i\) onto \(\mathbf{E}_i'\) ( \(1 \leqslant i \leqslant n\) ). Then \((f_1, \ldots, f_n)\) is said to be an isomorphism of the sets \(\mathbf{E}_1, \ldots, \mathbf{E}_n\) , endowed with the structure \(\mathbf{U}\) , onto the sets \(\mathbf{E}_1', \ldots, \mathbf{E}_n'\) , endowed with the structure \(\mathbf{U}'\) , if we have (in \(\mathfrak{C}'\) )

\[
\langle f _ {1}, \dots , f _ {n}, \operatorname{Id} _ {1}, \dots , \operatorname{Id} _ {m} \rangle^ {\mathrm{s}} (\mathrm{U}) = \mathrm{U} ^ {\prime}\tag{4}
\]

where \(Id_{h}\) denotes the identity mapping of \(A_{h}\) onto itself ( \(1 \leqslant h \leqslant m\) ). ¶ Let \(f_{i}^{\prime}\) be the inverse of the bijection \(f_{i}\) ( \(1 \leqslant i \leqslant n\) ). It follows immediately from (4) and the criterion CST3 (no. 2) that we have

\[
\langle f _ {1} ^ {\prime}, \dots , f _ {n} ^ {\prime}, \operatorname{Id} _ {1}, \dots , \operatorname{Id} _ {m} \rangle^ {\mathrm{s}} (\mathbf {U} ^ {\prime}) = \mathbf {U}
\]

and consequently that \((f_1', \ldots, f_n')\) is an isomorphism of \(\mathbf{E}_1', \ldots, \mathbf{E}_n'\) , endowed with \(\mathbf{U}'\) , onto \(\mathbf{E}_1, \ldots, \mathbf{E}_n\) , endowed with \(\mathbf{U}\) . The isomorphisms \((f_1, \ldots, f_n)\) and \((f_1', \ldots, f_n')\) are said to be inverses of each other.

\(E_{1}^{\prime}, \ldots, E_{n}^{\prime}\) , endowed with \(U^{\prime}\) , are said to be isomorphic to \(E_{1}, \ldots, E_{n}\) , endowed with U, if there exists an isomorphism of \(E_{1}, \ldots, E_{n}\) onto \(E_{1}^{\prime}, \ldots, E_{n}^{\prime}\) ; in this case the structures U and \(U^{\prime}\) are said to be isomorphic. The above definitions, together with CST1, imply the following criterion:

CST4. Let U, \(U'\) , \(U''\) be three structures of the same species \(\Sigma\) on the principal base sets \(E_{1}, \ldots, E_{n}, E_{1}', \ldots, E_{n}', E_{1}'', \ldots, E_{n}''\) , respectively. Let \(f_{i}\) be a bijection of \(E_{i}\) onto \(E_{i}'\) , and let \(g_{i}\) be a bijection of \(E_{i}'\) onto \(E_{i}''\) ( \(1 \leqslant i \leqslant n\) ). If \((f_{1}, \ldots, f_{n})\) and \((g_{1}, \ldots, g_{n})\) are isomorphisms, then so is \((g_{1} \circ f_{1}, \ldots, g_{n} \circ f_{n})\) .

An isomorphism of \(E_{1}, \ldots, E_{n}\) onto \(E_{1}, \ldots, E_{n}\) (with respect to the same structure) is called an automorphism of \(E_{1}, \ldots, E_{n}\) . The composition of two automorphisms of \(E_{1}, \ldots, E_{n}\) is an automorphism, and so is the inverse of an automorphism, * so that the automorphisms of \(E_{1}, \ldots, E_{n}\) form a group. *

Remark. By abuse of language, if \(f_{i}\) is any bijection of \(\mathbf{E}_i\) onto \(\mathbf{E}_i'\) ( \(1 \leqslant i \leqslant n\) ), \((f_1, \ldots, f_n)\) is said to be an isomorphism of \(\mathbf{E}_1, \ldots, \mathbf{E}_n\) onto \(\mathbf{E}_1', \ldots, \mathbf{E}_n'\) with respect to the species of structure of a set (no. 4, Remark 3).

CST5. In a theory \(\mathfrak{G}'\) which is stronger than \(\mathfrak{G}\) , let U be a structure of species \(\Sigma\) on \(\mathbf{E}_1, \ldots, \mathbf{E}_n\) , and let \(f_i\) be a bijection of \(\mathbf{E}_i\) onto a set \(\mathbf{E}_i'\) ( \(1 \leqslant i \leqslant n\) ). Then there exists a unique structure of species \(\Sigma\) on \(\mathbf{E}_1', \ldots, \mathbf{E}_n'\) such that \((f_1, \ldots, f_n)\) is an isomorphism of \(\mathbf{E}_1, \ldots, \mathbf{E}_n\) onto \(\mathbf{E}_1', \ldots, \mathbf{E}_n'\) .

For this structure, if it exists, can only be the term \(U'\) defined by the relation (4); it remains to be verified that this term is indeed a structure of species \(\Sigma\) , i.e., that the relation \(R\{E_{1}', \ldots, E_{n}', U'\}\) is true in \(C'\) . But this follows from the fact that \(R\{x_{1}, \ldots, x_{n}, s\}\) is transportable, for

\(R\{E_{1}^{\prime},\ldots,E_{n}^{\prime},U^{\prime}\}\) is equivalent in \(C'\) to the relation \(R\{E_{1},\ldots,E_{n},U\}\) (no. 3), which is true in \(C'\) by hypothesis.

The structure \(U'\) is said to be obtained by transporting the structure U to the sets \(E_{1}^{\prime}, \ldots, E_{n}^{\prime}\) by means of the bijective mappings \(f_{1}, \ldots, f_{n}\) . Thus two structures of the same species are isomorphic if and only if each is obtained from the other by transport of structure.

It may happen that any two structures of species \(\Sigma\) are necessarily isomorphic; the species of structure \(\Sigma\) is then said be univalent. * This is the case for the structure of an infinite cyclic group (isomorphic to \(\mathbf{Z}\) ), the structure of a prime field of characteristic zero (isomorphic to \(\mathbf{Q}\) ), the structure of a complete Archimedean ordered field (isomorphic to \(\mathbf{R}\) ), the structure of an algebraically closed, connected, locally compact topological field (isomorphic to \(\mathbf{C}\) ), and of the structure of a non-commutative, connected, locally compact topological division ring (isomorphic to \(\mathbf{H}\) , the division ring of quaternions). For some of these species of structure, for example that of a prime field of characteristic zero, or that of a complete Archimedean ordered field, there is not even any automorphism other than the identity mapping; but there do exist such automorphisms for the other examples given above (for example, the symmetry \(x \to -x\) in \(\mathbf{Z}\) ). *

It will be observed that the above species of structure are essentially those which are the basis of classical mathematics. On the other hand, * the species of group structures, the species of structures of an ordered set, and the species of topological structures, are not univalent. *

